\documentclass[11pt]{article}

\newcommand{\cnum}{Math 245B}
\newcommand{\ced}{Winter 2026}
\newcommand{\ctitle}[4]{\title{\vspace{-0.5in}\cnum, \ced\\Assignment #1: #2}\author{\vspace{-0.35in}\\#3, #4}}
\usepackage{enumitem}
\usepackage[usenames,dvipsnames,svgnames,table,hyperref]{xcolor}
\usepackage{amsmath, amsfonts}
\usepackage{graphicx} % Required for inserting images
\usepackage{float}
\usepackage{listings}
\usepackage{xcolor}
\usepackage{hyperref}
\usepackage{comment}

\renewcommand*{\theenumi}{\alph{enumi}}
\renewcommand*\labelenumi{(\theenumi)}
\renewcommand*{\theenumii}{\roman{enumii}}
\renewcommand*\labelenumii{\theenumii.}

\definecolor{codegreen}{rgb}{0,0.6,0}
\definecolor{codegray}{rgb}{0.5,0.5,0.5}
\definecolor{codepurple}{rgb}{0.58,0,0.82}
\definecolor{backcolour}{rgb}{0.95,0.95,0.92}
\definecolor{header}{HTML}{205459}
\definecolor{solution}{HTML}{204d52}

\newcommand{\solution}[1]{{{\textcolor{header}{Solution:} \textcolor{solution}{#1}}}}
\setlist[enumerate]{label=(\alph*)}

\lstdefinestyle{mystyle}{
    backgroundcolor=\color{backcolour},
    commentstyle=\color{codegreen},
    keywordstyle=\color{magenta},
    numberstyle=\tiny\color{codegray},
    stringstyle=\color{codepurple},
    basicstyle=\ttfamily\footnotesize,
    breakatwhitespace=false,
    breaklines=true,
    captionpos=b,
    keepspaces=true,
    numbers=left,
    numbersep=5pt,
    showspaces=false,
    showstringspaces=false,
    showtabs=false,
    tabsize=2
}


\begin{document}
\ctitle{\#1}{}{Asher Christian}{006-150-286}
\date{}
\maketitle

\section{Reminder}
Yesterday's lecture followed Marek Biskup analysis textbook
$X$ is a non empty set. \\
Exercise 1.2: Let $S \subset 2^{X}$ be a semi-algebra and let $ \mathcal{A}$ be the
collection of finite unions of elements of  $S$ so that by exercise  $1.1$
$ \mathcal{A}$ is an algebra. Show that every finitely additive function $\mu_0$ on $S$
extends uniquely to a finitely additive function $\mu$ on  $ \mathcal{A}$ and $\mu$
is unique.\\
Exercise 1.4 + 1.6: Further assume in Exercise 1.2 that $\mu$ is countablly subadditive 
show that $\mu$ is also countably subadditive.
\section{Definition}
\emph{
    Let $X$ be a metric space and let $B(X)$ be the  $\sigma$-algebra generated by the open balls
    in $X$. 
    \begin{enumerate}
        \item
            We call  $B(X)$ the borel  $\sigma$-algebra on $X$. Any element of $B(X)$ is called a
            Borel set.
        \item
            If $\mu : 2^{X} \rightarrow [0,+\infty]$ is an outer measure ($\mu(\emptyset) = 0$ and $\mu $ is countably sub-additive)
            we say that $B \subset X$ is $ \mu $-measurable if for all $A \subset X$ we have that
            \[
                \mu(A) \ge \mu(A \setminus B) + \mu(A \cap B) \iff \text{equality}
            \] 
            We denote the sets of $ \mu $-measurable subsets by $\Sigma(\mu)$
        \item We call $\mu$ a Borel measure if $B(X) \subset \Sigma(\mu)$
    \end{enumerate}
}
Given $ \mu $ borel I can define
\[
    F(x) = \mu(-\infty,x] \qquad G(x) = \mu(0,x]
\] 
\section{Theorem 3}
\emph{
    Let $F: \mathbb{R} \rightarrow \mathbb{R}$  be a monotone, nondecreasing, right continuous function and let $\mathcal{S}$ be 
    the collection of sets  $(a,b]$ where  $-\infty \le a \le b \le +\infty$. Define
    \[
        \mu_0((a,b])=  F(b) - F(a) \qquad \text{if $(a,b] \in \mathcal{S}$}
    \] 
    Then there exists a unique measure $\mu_F $ on $B( \mathbb{R})$ which is an extension of $\mu_0$
}\\
Reformulation of theorem: Let $F : \mathbb{R} \rightarrow \mathbb{R}$ be a monotone, nondecreasing, right continuous function. Then there exists
a unique borel measure $\mu_F$ on $ \mathbb{R}$ such that, $\mu((a,b]) = F(b) - F(a)$ if $-\infty \le a \le b \le +\infty$
\\
Proof:
Let $\mathcal{S}$ be defined as above. and Let  $ \mathcal{A}$ be the collection of finite unions of elements of $\mathcal{S}$, By exercise $1.1$ 
$ \mathcal{S}$ is a semi-algebra, by Lemma 2(stated monday) $\mu_0$ defined on $\mathcal{S}$ by  $\mu_0((a,b]) = F(b) - F(a)$ if $(a,b] \in \mathcal{S}$
is additive and countably additive on $ \mathcal{A}$. by Exericse 1.4, $ \mathcal{A}$ is an algebra on $ \mathbb{R}$. By Exercise 1.2 $\mu_0$ extends uniquely to an additive 
measure $ \mu $ on $A$. By Exercise 1.4 and 1.6,  $\mu$ is countably subadditive. By Exercise 1.3, $\mu$ extends to an outer measure on $ \mathbb{R}$. One checks that $\mu$ is a measure on $B( \mathbb{R})$.

\section{Definition}
\emph{
    A measure on $ \mathbb{R}$ is called Lebesgue-Stieltjes if it is  of the form $\mu_F$ as in
    theorem 3.
}

\section{Definition}
\emph{
    A measure on $ \mathbb{R}^{d}$ is called lebesgue-stieltjes if it is a Borel measure which is
    finite on bounded sets.
}

Hilbert space $H$:
\begin{enumerate}
    \item $a,b \in H$
    \item $\exists \lambda \in \mathbb{R}, b = \lambda a$
    \item $a \perp b$  $\langle a , b \rangle = 0$
    \item
        Fact $\exists a_b^{||}, a_b^{\perp} \in H$ such that $b = a_b^{||} + a_b^{\perp}$ and $a_b^{||} || b $ $a^{\perp}_b \perp b$
\end{enumerate}
Set  $X$:
\begin{enumerate}
    \item $\mu,\nu$ measure on $X$
    \item  $\nu << \mu$ (  $\nu$ is absolutely continous with respect to $ \mu $ )\\
        $\nu(A) = \int_{A}^{}f d \mu$ \\
        \[
            f(x) = \lim_{x\to 0}\frac{\nu(B(x,r))}{\mu(B(x,r))} = D_{\mu}\nu
        \] 
    \item $\mu \perp \nu$ ($u$ and $\nu$ are mutually orthongonal)
    \item $\nu = \nu_{\text{abs}} + \nu_{\text{perp}}$  $\nu_{\text{abs}} << \mu$, $\nu_{\text{perp}} \perp \mu$
\end{enumerate}

\section{Definition}
\emph{
    Let $\mu$ and $\nu$ be Borel measure on a metric space $ X$
    \begin{enumerate}
        \item We say that $ \nu $ is absolutely continuous with respect to $\mu$ and we write
            $\nu << \mu$ if whenver $A \subset X$ and $\mu(A) = 0$ then $\nu(A) = 0$ (Evans - Gariepy Fine properties of functions)
        \item We say that $\mu$ and $\nu$ are mutually orthogonal (perpendicular) if there exists  $B \subset \mathbb{R}$ Borel such that $\mu(B) = 0. \nu(X \setminus B) = 0$
    \end{enumerate}
}
Example $ \mu$ is length of $ \mathbb{R}$ and
\[
    \delta_{x_0}(A) = \nu(A) = 
    \begin{cases}
        0 & x \not\in A \\
        1 & x \in A
    \end{cases}
\] 
then if $B =  \mathbb{R}\setminus \{x_0\}$ we can see that $\mu \perp \nu$\
\section{X is a metric space}
Reminder: Let $ \mu $ be an outer measure on $X$, we say that  $\mu$ is Borel regular if every borel set is $\mu$-measurable and 
for every $A \subset X$ there exists $ B \subset X$ Borel such that $A \subset B$ and
\[
    \mu[A] = \mu[B]
\] 
we say that $ \mu $ is a Radon measure if $\mu$ is Borel regular and $\mu[K] < +\infty$ for all $ K \subset  X $ compact. Two covering lemmas (Evans-Gariepy)

\section{Theorem Besi courtch covering's lemma}
\emph{
    Closed ball: $B_r(x) := \{y \in X: d(x,y) \le r\}$  $r > 0$ non degenerate
}
\emph{
    Given $d \in \mathbb{N}$ there exists $ N_d > 1$ natural number such that the following holds.
    If  $F$ is a collection of closed non degenerate balls in $ \mathbb{R}^{d}$ such that
     \[
          R := \sup_{B \in F} \text{diam}(B) < +\infty
    \] 
    then there are countable families $g_1, \dots, g_{N_d}$ of balls from $F$ such that.
    each  $g_i$ is made of disjoint balls and 
     \[
         A \subset \bigcup_{i=1}^{N_d}g_i
    \] 
    where $A$ is the collection of the centers of the balls in $ F$
}
\[
    x_0 \in A \implies x_0 \in B_r(a)
\] 
where $ B_r(a)$ intersects  at most $N_d - 1$ balls of  $\bigcup_{i=1}^{N_d}g_i$

\section{Theorem [Vitali's covering Lemma: works in separable metric space]}
\emph{
    Let $F$ be a collection of non degenerate closed balls such that
    \[
        R := \sup_{B \in F}\text{diam}(B) < +\infty
    \] 
    then there exists a countable family $G$ of balls in $F$ such that 
    the balls in $G$ are disjoint and
     \[
         \bigcup_{B \in F} B \subset \bigcup_{B \in G} \hat B
    \] 
    if $B = B_r(a)$ then $\hat B = B_{5r}(a)$
}
Counter example when $\sup_{B} \text{diam}(B) = +\infty$ use balls centered at zero going to infinity.

\section{Definition of cover}
\emph{
    A collection  $F$ of balls in  $X$ is called  a cover of $A \subset X$ if 
    \[
        A \subset \bigcup_{B \in F} B
    \] 
    If we further assume that for every $x \in A$
     \[
         \inf_{B \in F}\left\{ \text{diam}(B) : B \in F \;\; x \in B \right\} = 0
    \] 
    we call the cover fine.
}

\section{corollary 1}
\emph{
    Assume that $F$ si a family of closed nondegenerate balls in $ \mathbb{R}^{d}$ and $A \subset  \mathbb{R}^{d}$ is such that $F$ is a fine cover
    of  $A$. If 
    \[
        \sup_{B \in F} \text{diam})B < +\infty
    \] 
    then there exists $G$ a collection of disjoint balls of  $F$ such that
    whenever I choose   $B_1, \dots, B_m \in F$ then
    \[
        A \setminus (B_1 \cup \dots \cup B_m) \subset \bigcup_{B \in G \setminus\{B_1,\dots,B_m\}} \hat B
    \] \\
}
Proof: The family $G$ in the Vitali's covering lemma satisfies the following property:\\
if $B \in F, B' \in G$ and  $B \cap B' \ne \emptyset$ then  $B \subset \hat B'$.\\
Let $B_1, \dots, B_m \in F$
Case 1 $A \subset B_1 \cup \dots \cup B_m$ then $A \setminus (b_1 \cup \dots \cup B_m) = \emptyset$ and there is nothing to prove\\
Case 2 $A \not\subset B_1 \cup \dots \cup B_m$ let $x \in A \setminus (B_1 \cup \dots \cup B_m)$ $x \in \bigcup_{B \in F} B$  $x \not\in B_1 \cup \dots \cup B_m$ which is a closed set so we can find a 
radius so small so that $B_r(x) \cap (B_1 \cup \dots \cup B_m) = \emptyset$ then there exists another $B \in F$ such that  $B \subset B_r(x)$. Choose $B' \in G$ such that  $B \cap B' \ne \emptyset$ we have that
 $x \in B \subset \hat B'$ hence $x \in \hat B'$

\section{Corollary 2}
\emph{
    Let $O \subset \mathbb{R}^{d}$ be an open set and let $\delta > 0$. Then there exists
    a collection $G$ of disjoint balls non degenerate closed contained in  $O$ such that
     \[
         \sup_{B \in G} \text{diam}(B) \le \delta
    \] 
    and
    \[
        \mathcal{L}^{d}(O \setminus \bigcup_{B \in G} B) = 0
    \] 
}



 

\end{document}
